\chapter{Derivation of the StageBridge Training Objective and Estimand} \label{chap:appendix-b}

%% Appendix B: derivation of the model. (1) flow matching -> why L_CFM recovers the OT displacement
%% field; (2) the structured UDE decomposition + why the three-stage schedule identifies self vs
%% context; (3) why the self-vs-full endpoint difference is the within-receiver context estimand R_cond;
%% (4) the full objective. Matches crrt/ude/train.py (CFM + 0.25 endpoint + 0.01 residual + 1e-4 sparse
%% + anchor) and crrt/ude/estimand.py. VOICE: numbered equations, every symbol defined, plain-word claim.

This appendix derives the StageBridge training objective and shows why the context residual
$R_{\mathrm{cond}}$ of Chapter~\ref{chap:methods} (Equation~\ref{eq:methods-rcond-identity}) is a valid within-receiver
measure of the effect of local context on progression. The argument has three parts: flow matching
against an optimal-transport coupling recovers the displacement field between the source and target
regulatory distributions; the three-stage schedule identifies the self and context components of that
field; and the difference between the full and self endpoints, integrated from a shared initial state,
isolates the context contribution.

\section{Flow matching against a transport coupling}

Let $\mu_0$ and $\mu_1$ be the source-stage and target-stage distributions of regulatory states, and let
$\Pi^{\star}$ be the entropic optimal-transport coupling between samples of them
(Appendix~\ref{chap:appendix-a}). For a coupled pair $(z_i, z_j) \sim \Pi^{\star}$, define the linear
interpolant and its constant velocity,
\begin{equation}\label{eq:app-interp}
z_\tau \;=\; (1-\tau)\, z_i + \tau\, z_j ,
\qquad
\frac{d}{d\tau} z_\tau \;=\; z_j - z_i \;\equiv\; u_{ij} ,
\end{equation}
where $z_\tau$ is the state at pseudo-time $\tau \in [0,1]$ along the straight path from $z_i$ to $z_j$
and $u_{ij}$ is the displacement velocity of that path. Conditional flow matching \autocite{Lipman2022-ns} trains a
velocity field $v(z, w, \tau)$ to regress onto this conditional velocity,
\begin{equation}\label{eq:app-cfm}
\mathcal{L}_{\mathrm{CFM}} \;=\;
\mathbb{E}_{(i,j)\sim\Pi^{\star},\ \tau\sim U[0,1]}
\left\lVert\, v\!\left(z_\tau, w_i, \tau\right) - u_{ij} \,\right\rVert_2^2 ,
\end{equation}
where the expectation is over coupled pairs and over pseudo-time. The minimizer of
Equation~\ref{eq:app-cfm} at a point $(z, \tau)$ is the conditional mean of the displacement velocity
over all coupled pairs whose interpolant passes through that point,
\begin{equation}\label{eq:app-marginal}
v^{\star}(z, \tau) \;=\; \mathbb{E}_{(i,j)\sim\Pi^{\star}}\!\left[\, u_{ij} \ \middle|\ z_\tau = z \,\right] ,
\end{equation}
which is the marginal velocity field whose flow transports $\mu_0$ to $\mu_1$ along the transport plan
$\Pi^{\star}$. In words: regressing the field onto straight-line displacements of transport-matched pairs
recovers the average velocity that carries the source distribution onto the target distribution, so the
learned $v$ integrates source cells to the target-stage regulatory distribution. Because $\Pi^{\star}$
is computed only from regulatory-state distance within a coarse niche stratum
(Equation~\ref{eq:methods-entropic-coupling}), the supervision carries no fine-context information, and the context
dependence of $v$ is learned rather than imposed by the pairing.

\section{The structured decomposition and its estimation}
\label{sec:app-b-decomposition}

StageBridge does not fit an unstructured $v$. It writes the field as a receiver-intrinsic term plus a
context term,
\begin{equation}\label{eq:app-decomp}
v(z, w, \tau) \;=\; \underbrace{v_{\mathrm{self}}(z, \tau)}_{\text{no }w}
\;+\; \underbrace{v_{\mathrm{context}}(z, w, \tau)}_{\text{vanishes when context is inert}} ,
\end{equation}
with $v_{\mathrm{self}}$ and $v_{\mathrm{context}}$ as defined in
Equations~\ref{eq:methods-self-field} and~\ref{eq:methods-context-field}. This decomposition is identifiable only if the two terms
are estimated so that $v_{\mathrm{self}}$ captures the progression that occurs independent of context
and $v_{\mathrm{context}}$ captures only the residual that context adds. A single joint fit does not
guarantee this: the flexible context term could absorb intrinsic progression, or the self term could
absorb context-driven variation, and the split would be arbitrary. The three-stage schedule resolves
this.

\emph{Stage A} fits $v_{\mathrm{self}}$ alone against the coupling velocity, minimizing
Equation~\ref{eq:app-cfm} in self mode (the field evaluated with $v_{\mathrm{context}} \equiv 0$). Its
minimizer is the best context-free explanation of the transport,
\begin{equation}\label{eq:app-stageA}
v_{\mathrm{self}}^{A} \;=\; \argmin_{v_{\mathrm{self}}}\ \mathbb{E}
\left\lVert\, v_{\mathrm{self}}(z_\tau, \tau) - u_{ij} \,\right\rVert_2^2 .
\end{equation}
\emph{Stage B} freezes $v_{\mathrm{self}} = v_{\mathrm{self}}^{A}$ and fits the context term to the
part of the displacement the frozen self field does not explain,
\begin{equation}\label{eq:app-stageB}
v_{\mathrm{context}}^{B} \;=\; \argmin_{v_{\mathrm{context}}}\ \mathbb{E}
\left\lVert\, v_{\mathrm{context}}(z_\tau, w_i, \tau) - \bigl(u_{ij} - v_{\mathrm{self}}^{A}(z_\tau, \tau)\bigr) \,\right\rVert_2^2 ,
\end{equation}
so the context term regresses onto the \emph{residual} velocity $u_{ij} - v_{\mathrm{self}}^{A}$ by
construction. This is what ties $v_{\mathrm{context}}$ to what is left over after intrinsic progression,
which is the operational meaning of ``context residual.'' \emph{Stage C} refines both terms jointly at
a reduced learning rate under an anchor penalty that keeps the self field close to its Stage-A value,
\begin{equation}\label{eq:app-anchor}
\mathcal{L}_{\mathrm{anchor}} \;=\; \lambda_a\, \frac{1}{P}\sum_{k=1}^{P}
\bigl\lVert \theta_{s,k} - \theta_{s,k}^{A} \bigr\rVert_2^2 ,
\end{equation}
where $\theta_{s,k}$ is the $k$-th self-field parameter, $\theta_{s,k}^{A}$ its Stage-A value, $P$ the
total self-field parameter count, and $\lambda_a$ the anchor weight. Normalizing by $P$ makes
$\lambda_a$ independent of model size. The anchor lets the joint step correct small mismatches at the
interface of the two terms while preventing Stage C from silently relearning the self field as context,
so the identification established in Stages A and B is preserved.

\section{The context residual as a within-receiver contrast}

The trained field defines, for a receiver with initial state $z_0$ and context $w$, two endpoints
obtained by integrating from the \emph{same} $z_0$,
\begin{align}
z_{\mathrm{full}}(1) \;&=\; z_0 + \int_0^1 \bigl[\, v_{\mathrm{self}}(z(\tau), \tau) + v_{\mathrm{context}}(z(\tau), w, \tau) \,\bigr]\, d\tau , \label{eq:app-full}\\
z_{\mathrm{self}}(1) \;&=\; z_0 + \int_0^1 v_{\mathrm{self}}(z(\tau), \tau)\, d\tau , \label{eq:app-selfend}
\end{align}
where the full trajectory includes the context modulation and the self trajectory omits it, and both
share the identical self parameters. The context residual is their difference under the semantic
projection $\phi$,
\begin{equation}\label{eq:app-rcond}
R_{\mathrm{cond}}(z_0, w) \;=\; \phi\!\left[z_{\mathrm{full}}(1)\right] - \phi\!\left[z_{\mathrm{self}}(1)\right]
\;\xrightarrow{\ \phi=\mathrm{id}\ }\;
z_{\mathrm{full}}(1) - z_{\mathrm{self}}(1) .
\end{equation}
Because the two integrals in Equations~\ref{eq:app-full}--\ref{eq:app-selfend} differ only by the
presence of $v_{\mathrm{context}}$ and are driven from a common $z_0$ with shared self parameters,
$R_{\mathrm{cond}}$ is a \emph{within-receiver} contrast: it holds the receiver's own state and intrinsic
dynamics fixed and reports only the displacement attributable to its context. It is not a comparison
between two separately fit models, which is the property that makes it interpretable as the effect of
context rather than an artifact of two independent fits. When context is inert
($v_{\mathrm{context}} \to 0$) the two trajectories coincide and $R_{\mathrm{cond}} \to 0$; coordinate
$d$ of $R_{\mathrm{cond}}$ is the context effect on regulatory program $d$. Fixing the context at the
observed source niche $w(\tau) = w_i$ throughout both integrals makes $R_{\mathrm{cond}}$ the effect of
the receiver's \emph{observed} ecology, consistent with the integrator, which evaluates a single fixed
context at every sub-step.

\section{The full training objective}

The complete per-step objective adds three regularizers to the flow-matching term. The endpoint term
matches the integrated source population to the target population in transport distance, the residual
term keeps the context velocity small so context acts as a correction rather than a driver, and the
sparsity term encourages sparse context effects,
\begin{equation}\label{eq:app-objective}
\mathcal{L} \;=\; \mathcal{L}_{\mathrm{CFM}}
\;+\; \lambda_{e}\, d_{\mathrm{Sinkhorn}}\!\bigl(z_{\mathrm{int}}(1),\, \mu_1\bigr)
\;+\; \lambda_{r}\, \mathbb{E}\lVert v_{\mathrm{context}} \rVert_2^2
\;+\; \lambda_{s}\, \lVert \Theta_{c} \rVert_1
\;+\; \mathbf{1}[\text{Stage C}]\,\mathcal{L}_{\mathrm{anchor}} ,
\end{equation}
where $z_{\mathrm{int}}(1)$ is the integrated endpoint of the source batch, $d_{\mathrm{Sinkhorn}}$ is
the entropic transport distance of Appendix~\ref{chap:appendix-a}, $\mu_1$ is the target-stage
population, $\Theta_{c}$ are the context-field parameters, $\mathbf{1}[\cdot]$ is the indicator that the
anchor is active only in Stage C, and $(\lambda_e, \lambda_r, \lambda_s, \lambda_a) = (0.25,\, 0.01,\, 10^{-4},\, 0.1)$
are the fixed loss weights, matching Equation~\ref{eq:methods-training-loss-weights}. The flow-matching term supplies the primary supervision of the velocity
field; the endpoint term corrects accumulated integration error at $\tau=1$; and the residual and
sparsity terms enforce that the context contribution is a small, sparse residual on top of the intrinsic
progression, matching the interpretation of $R_{\mathrm{cond}}$ as a residual effect.

\section{Three-stage training procedure in full}
\label{sec:app-b-training-detail}

Section~\ref{sec:methods-training} states the three-stage schedule
(Table~\ref{tab:training-schedule}), the composite Stage~C objective
(Equation~\ref{appb-eq:methods-stage-c-objective}), and the evaluated loss
coefficients (Equation~\ref{appb-eq:methods-training-loss-weights}); the
identification argument for the schedule is given in
Section~\ref{sec:app-b-decomposition} above. This section gives the remaining
implementation in full: the construction of each flow-matching training
example, the per-stage objective forms, each auxiliary term, and the numerical
integration scheme. It is included for reproduction of the fitted models.
Notation follows Chapter~\ref{chap:methods} throughout and no quantity is
redefined here.

The conditional-stratified transport procedure identifies probabilistic pairs
of source- and target-stage receivers. The StageBridge differential model
defines a velocity at any regulatory state located between those two stage
distributions. Conditional flow matching connects these components by training
the model velocity to follow the direction implied by each transport-sampled
source--target pair \autocite{Lipman2022-ns,Bunne2022-go}.

Flow matching converts endpoint observations into local velocity-supervision
examples. For a paired source state and target state, a point is sampled along
the straight line connecting them. The model is evaluated at that intermediate
point and trained to predict the displacement from the source state to the
target state. Repeating this procedure across transport-sampled pairs and
progression coordinates produces a continuous vector field whose integrated
trajectories connect the source and target populations.

Training was performed independently within each donor-held-out
cross-validation fold. Receiver states, ecological strata, transport
couplings, and model parameters used for fitting were constructed exclusively
from the training donors assigned to that fold.

\subsection{Flow-matching training examples}
\label{appb-sec:methods-cfm-paths}

Each optimization step began with the conditional-stratified transport plan
described in Section~\ref{sec:methods-entropic-ot}. From this plan,
\(B=64\) source--target pairs were sampled. Let
\(\nu\in\{1,\ldots,B\}\) index one of these sampled pairs, and let
\((i_\nu,j_\nu)\) denote its source and target receiver indices. Thus,
\(i_\nu\in\mathcal{S}\) identifies a source-stage receiver and
\(j_\nu\in\mathcal{T}\) identifies its transport-sampled target-stage
receiver.

For each pair, a progression coordinate was sampled independently:

\begin{equation}
\tau_\nu
\sim
\operatorname{Uniform}(0,1).
\label{appb-eq:methods-cfm-time-sampling}
\end{equation}

Here, \(\operatorname{Uniform}(0,1)\) is the continuous uniform distribution
on the interval from \(0\) to \(1\). Every value in this interval therefore
has equal sampling density. The scalar \(\tau_\nu\) specifies the position at
which the model is evaluated along the progression path for pair \(\nu\).

The intermediate regulatory state was constructed by linear interpolation:

\begin{equation}
\widetilde{z}_\nu
=
\left(1-\tau_\nu\right)
z_{i_\nu}^{\mathrm{obs}}
+
\tau_\nu
z_{j_\nu}^{\mathrm{obs}}.
\label{appb-eq:methods-cfm-path}
\end{equation}

The vector
\(\widetilde{z}_\nu\in\mathbb{R}^{d_z}\) is the interpolated regulatory state
used as the input to the differential field. At \(\tau_\nu=0\), it equals the
observed source state \(z_{i_\nu}^{\mathrm{obs}}\). At \(\tau_\nu=1\), it
equals the observed target state \(z_{j_\nu}^{\mathrm{obs}}\). Intermediate
values form a straight path between the two states in the
\(d_z\)-dimensional regulatory space.

Because Equation~\ref{appb-eq:methods-cfm-path} is linear in \(\tau_\nu\), its
derivative with respect to the progression coordinate is constant:

\begin{equation}
u_\nu
=
\frac{d\widetilde{z}_\nu}{d\tau_\nu}
=
z_{j_\nu}^{\mathrm{obs}}
-
z_{i_\nu}^{\mathrm{obs}}.
\label{appb-eq:methods-cfm-target}
\end{equation}

The vector \(u_\nu\in\mathbb{R}^{d_z}\) is the target velocity for training
pair \(\nu\). Each coordinate of \(u_\nu\) gives the change in one pathway or
transcription-factor activity between the paired source and target states.
The target is identical to the transport-paired displacement
\(\Delta z_{i_\nu j_\nu}\) defined in
Equation~\ref{eq:methods-pair-displacement}.

The ecological input was held at the observed source value throughout the
training path:

\begin{equation}
\widetilde{w}_\nu(\tau)
=
w_{i_\nu},
\qquad
\tau\in[0,1].
\label{appb-eq:methods-cfm-source-context}
\end{equation}

Here, \(\widetilde{w}_\nu(\tau)\) denotes the context supplied to the model
for training pair \(\nu\) at progression coordinate \(\tau\).
Equation~\ref{appb-eq:methods-cfm-source-context} states that this input remains
equal to the observed source context \(w_{i_\nu}\) across the complete
interval. The same source-fixed convention was used for pointwise training,
endpoint integration, and calculation of the conditional context residual.

\subsection{Velocity-matching loss}
\label{appb-sec:methods-cfm-loss}

For a predicted velocity \(\widehat{u}_\nu\in\mathbb{R}^{d_z}\), the
flow-matching error for pair \(\nu\) is the squared Euclidean distance between
the predicted and target velocities. Averaging this error across the \(B\)
sampled pairs and \(d_z\) regulatory coordinates gives

\begin{equation}
\mathcal{L}_{\mathrm{CFM}}
=
\frac{1}{B d_z}
\sum_{\nu=1}^{B}
\left\|
\widehat{u}_\nu-u_\nu
\right\|_2^2.
\label{appb-eq:methods-general-cfm-loss}
\end{equation}

Here, \(\mathcal{L}_{\mathrm{CFM}}\) is the conditional flow-matching loss,
\(\|\cdot\|_2\) is the Euclidean norm, and
\(\widehat{u}_\nu\) is the velocity predicted by the field used in the current
training stage. Division by \(B d_z\) expresses the loss as the mean squared
error per pair and regulatory coordinate. Minimizing this quantity aligns the
modeled local velocity with the displacement implied by the
transport-sampled pair.

StageBridge used three sequential training stages so that the receiver-state
and context-conditioned contributions were established separately before
joint refinement. Let \(\theta_s\) denote all trainable parameters of
\(v_{\mathrm{self}}\), and let \(\theta_c\) denote all trainable parameters of
\(v_{\mathrm{context}}\).

\subsection{Receiver-state field fitting}
\label{appb-sec:methods-training-stage-a}

Stage A fitted the receiver-state field. For pair \(\nu\), the predicted
velocity was

\begin{equation}
\widehat{u}_\nu^{(\mathrm{A})}
=
v_{\mathrm{self}}
\left(
\widetilde{z}_\nu,
\tau_\nu;
\theta_s
\right),
\label{appb-eq:methods-stage-a-prediction}
\end{equation}

where the superscript \(\mathrm{A}\) identifies Stage A and the semicolon
separates the field inputs from its trainable parameters. The self field
receives the interpolated receiver state and progression coordinate and
receives no ecological-context vector.

Substitution into
Equation~\ref{appb-eq:methods-general-cfm-loss} gives the Stage A flow-matching
loss:

\begin{equation}
\mathcal{L}_{\mathrm{CFM}}^{(\mathrm{A})}
=
\frac{1}{B d_z}
\sum_{\nu=1}^{B}
\left\|
v_{\mathrm{self}}
\left(
\widetilde{z}_\nu,
\tau_\nu;
\theta_s
\right)
-
u_\nu
\right\|_2^2.
\label{appb-eq:methods-stage-a-cfm}
\end{equation}

Stage A therefore estimated the portion of transport-aligned velocity
predictable from epithelial regulatory state and relative progression
position. Its training pairs remained restricted to the coarse ecological
common-support domain established by the conditional coupling.

At the end of Stage A, the fitted self-field parameters were stored as
\(\theta_s^{(\mathrm{A})}\). The resulting field was

\begin{equation}
\widehat{v}_{\mathrm{self}}^{(\mathrm{A})}(z,\tau)
=
v_{\mathrm{self}}
\left(
z,\tau;
\theta_s^{(\mathrm{A})}
\right).
\label{appb-eq:methods-stage-a-field}
\end{equation}

The notation
\(\widehat{v}_{\mathrm{self}}^{(\mathrm{A})}\) denotes the self field after
completion of Stage A. This stored field served as the fixed receiver-state
component in Stage B and as the reference solution for the Stage C anchor.

\subsection{Context-field fitting}
\label{appb-sec:methods-training-stage-b}

Stage B fixed the self-field parameters at
\(\theta_s^{(\mathrm{A})}\) and optimized the context-field parameters
\(\theta_c\). The full velocity predicted for pair \(\nu\) was

\begin{equation}
\widehat{u}_\nu^{(\mathrm{B})}
=
\widehat{v}_{\mathrm{self}}^{(\mathrm{A})}
\left(
\widetilde{z}_\nu,
\tau_\nu
\right)
+
v_{\mathrm{context}}
\left(
\widetilde{z}_\nu,
w_{i_\nu},
\tau_\nu;
\theta_c
\right).
\label{appb-eq:methods-stage-b-prediction}
\end{equation}

The first term is the frozen Stage A self-field prediction. The second term is
the context-associated velocity predicted from the interpolated state,
observed source context, and progression coordinate.

The Stage B flow-matching loss was

\begin{equation}
\mathcal{L}_{\mathrm{CFM}}^{(\mathrm{B})}
=
\frac{1}{B d_z}
\sum_{\nu=1}^{B}
\left\|
\widehat{u}_\nu^{(\mathrm{B})}
-
u_\nu
\right\|_2^2.
\label{appb-eq:methods-stage-b-cfm}
\end{equation}

Because the self field is fixed during this stage, minimizing
Equation~\ref{appb-eq:methods-stage-b-cfm} fits the context field to the velocity
remaining after the Stage A prediction. This remaining velocity can be written
explicitly as

\begin{equation}
u_\nu^{\mathrm{res}}
=
u_\nu
-
\widehat{v}_{\mathrm{self}}^{(\mathrm{A})}
\left(
\widetilde{z}_\nu,
\tau_\nu
\right),
\label{appb-eq:methods-stage-b-residual-target}
\end{equation}

where \(u_\nu^{\mathrm{res}}\in\mathbb{R}^{d_z}\) is the residual velocity
target for pair \(\nu\). Stage B therefore estimates the context-field
contribution required to supplement the fitted receiver-state field.

\subsection{Anchored joint refinement}
\label{appb-sec:methods-training-stage-c}

Stage C allowed the parameters of both fields to update jointly. The predicted
velocity was

\begin{equation}
\widehat{u}_\nu^{(\mathrm{C})}
=
v_{\mathrm{self}}
\left(
\widetilde{z}_\nu,
\tau_\nu;
\theta_s
\right)
+
v_{\mathrm{context}}
\left(
\widetilde{z}_\nu,
w_{i_\nu},
\tau_\nu;
\theta_c
\right).
\label{appb-eq:methods-stage-c-prediction}
\end{equation}

The Stage C flow-matching loss was

\begin{equation}
\mathcal{L}_{\mathrm{CFM}}^{(\mathrm{C})}
=
\frac{1}{B d_z}
\sum_{\nu=1}^{B}
\left\|
\widehat{u}_\nu^{(\mathrm{C})}
-
u_\nu
\right\|_2^2.
\label{appb-eq:methods-stage-c-cfm}
\end{equation}

Joint optimization permits coordinated adjustment of the two fields. An
anchor penalty limits displacement of the self-field parameters from their
Stage A values. Let \(P_s\) be the number of scalar parameters in
\(\theta_s\), and let \(\theta_{s,p}\) denote scalar self-field parameter
\(p\). The anchor penalty was

\begin{equation}
\mathcal{L}_{\mathrm{anchor}}
=
\frac{1}{P_s}
\sum_{p=1}^{P_s}
\left(
\theta_{s,p}
-
\theta_{s,p}^{(\mathrm{A})}
\right)^2.
\label{appb-eq:methods-stage-c-anchor}
\end{equation}

Here, \(\theta_{s,p}^{(\mathrm{A})}\) is the stored value of parameter \(p\)
at the end of Stage A. Dividing by \(P_s\) gives the mean squared parameter
change and prevents the nominal anchor strength from increasing automatically
with model size. The anchor retains the Stage A receiver-state field as the
reference decomposition while permitting limited joint refinement.

\subsection{Endpoint population loss}
\label{appb-sec:methods-endpoint-loss}

Flow matching supervises velocity at sampled intermediate states. StageBridge
also compared the population of integrated endpoints with the sampled
target-state population.

For training stage
\(\kappa\in\{\mathrm{A},\mathrm{B},\mathrm{C}\}\), let
\(\widehat{z}_{\nu,1}^{(\kappa)}\) denote the endpoint obtained by integrating
the field used in stage \(\kappa\) from source state
\(z_{i_\nu}^{\mathrm{obs}}\) at \(\tau=0\) to \(\tau=1\). The matrix of
predicted endpoints was

\begin{equation}
\widehat{Z}_1^{(\kappa)}
=
\begin{bmatrix}
\left(\widehat{z}_{1,1}^{(\kappa)}\right)^{\mathsf T}\\
\vdots\\
\left(\widehat{z}_{B,1}^{(\kappa)}\right)^{\mathsf T}
\end{bmatrix}
\in
\mathbb{R}^{B\times d_z}.
\label{appb-eq:methods-predicted-endpoint-matrix}
\end{equation}

Here, each row contains one predicted regulatory endpoint. Stage A endpoints
were generated with the self field. Stage B endpoints were generated with the
frozen self field plus the trainable context field. Stage C endpoints were
generated with the jointly refined full field.

The corresponding sampled target-state matrix was

\begin{equation}
Z_{\mathrm{tgt}}
=
\begin{bmatrix}
\left(z_{j_1}^{\mathrm{obs}}\right)^{\mathsf T}\\
\vdots\\
\left(z_{j_B}^{\mathrm{obs}}\right)^{\mathsf T}
\end{bmatrix}
\in
\mathbb{R}^{B\times d_z}.
\label{appb-eq:methods-training-target-matrix}
\end{equation}

The endpoint loss was

\begin{equation}
\mathcal{L}_{\mathrm{endpoint}}^{(\kappa)}
=
\operatorname{Sinkhorn}_{\varepsilon}
\left(
\widehat{Z}_1^{(\kappa)},
Z_{\mathrm{tgt}}
\right).
\label{appb-eq:methods-endpoint-loss}
\end{equation}

The function
\(\operatorname{Sinkhorn}_{\varepsilon}(X,Y)\) is the entropically
regularized optimal-transport discrepancy between the empirical point sets
\(X\) and \(Y\), using uniform weights and squared Euclidean distance in
regulatory space. It measures how closely the distribution of predicted
endpoints matches the distribution of sampled target states. This comparison
is population-level: it allows transport mass to be distributed across the
endpoint and target sets instead of requiring every integrated source to end
exactly at its sampled target partner. The endpoint calculation used
\(\varepsilon=0.05\) and \(80\) Sinkhorn iterations, matching the transport
configuration used for pair construction.

\subsection{Context-field regularization}
\label{appb-sec:methods-context-regularization}

Two penalties regulated the magnitude and parameterization of the context
field. The first penalized the mean squared context velocity:

\begin{equation}
\mathcal{L}_{\mathrm{residual}}
=
\frac{1}{B d_z}
\sum_{\nu=1}^{B}
\left\|
v_{\mathrm{context}}
\left(
\widetilde{z}_\nu,
w_{i_\nu},
\tau_\nu;
\theta_c
\right)
\right\|_2^2.
\label{appb-eq:methods-context-residual-penalty}
\end{equation}

This term favors a context field whose contribution remains small unless a
larger context-associated velocity improves agreement with the transport
targets. It was active during Stages B and C, when the full field was used.

The second penalty acted on the context-field parameters. Let \(P_c\) be the
number of scalar parameters in \(\theta_c\), and let \(\theta_{c,p}\) denote
context-field parameter \(p\). The sparsity penalty was

\begin{equation}
\mathcal{L}_{\mathrm{sparse}}
=
\frac{1}{P_c}
\sum_{p=1}^{P_c}
\left|
\theta_{c,p}
\right|.
\label{appb-eq:methods-context-sparsity}
\end{equation}

This is the mean absolute parameter magnitude, equivalent to a
parameter-normalized \(L_1\) penalty. It favors context-field parameterizations
containing more values near zero. The penalty was included throughout the
three-stage schedule. During Stage A, it regularized the inactive context
branch while the self field supplied the velocity prediction.

\subsection{Stage-specific objectives}
\label{appb-sec:methods-stage-objectives}

The complete objective differed across the three stages. Stage A minimized

\begin{equation}
\mathcal{L}^{(\mathrm{A})}
=
\mathcal{L}_{\mathrm{CFM}}^{(\mathrm{A})}
+
\lambda_{\mathrm{endpoint}}
\mathcal{L}_{\mathrm{endpoint}}^{(\mathrm{A})}
+
\lambda_{\mathrm{sparse}}
\mathcal{L}_{\mathrm{sparse}}.
\label{appb-eq:methods-stage-a-objective}
\end{equation}

Stage B minimized

\begin{equation}
\mathcal{L}^{(\mathrm{B})}
=
\mathcal{L}_{\mathrm{CFM}}^{(\mathrm{B})}
+
\lambda_{\mathrm{endpoint}}
\mathcal{L}_{\mathrm{endpoint}}^{(\mathrm{B})}
+
\lambda_{\mathrm{residual}}
\mathcal{L}_{\mathrm{residual}}
+
\lambda_{\mathrm{sparse}}
\mathcal{L}_{\mathrm{sparse}}.
\label{appb-eq:methods-stage-b-objective}
\end{equation}

Stage C minimized

\begin{equation}
\begin{aligned}
\mathcal{L}^{(\mathrm{C})}
={}&
\mathcal{L}_{\mathrm{CFM}}^{(\mathrm{C})}
+
\lambda_{\mathrm{endpoint}}
\mathcal{L}_{\mathrm{endpoint}}^{(\mathrm{C})}
\\
&+
\lambda_{\mathrm{residual}}
\mathcal{L}_{\mathrm{residual}}
+
\lambda_{\mathrm{sparse}}
\mathcal{L}_{\mathrm{sparse}}
+
\lambda_{\mathrm{anchor}}
\mathcal{L}_{\mathrm{anchor}}.
\end{aligned}
\label{appb-eq:methods-stage-c-objective}
\end{equation}

The symbols
\(\lambda_{\mathrm{endpoint}}\),
\(\lambda_{\mathrm{residual}}\),
\(\lambda_{\mathrm{sparse}}\), and
\(\lambda_{\mathrm{anchor}}\)
are nonnegative coefficients controlling the contribution of each auxiliary
loss. Their evaluated values were

\begin{equation}
\lambda_{\mathrm{endpoint}}=0.25,
\qquad
\lambda_{\mathrm{residual}}=0.01,
\qquad
\lambda_{\mathrm{sparse}}=10^{-4},
\qquad
\lambda_{\mathrm{anchor}}=0.1.
\label{appb-eq:methods-training-loss-weights}
\end{equation}

The flow-matching term has implicit weight \(1\) and defines the primary
velocity-fitting objective. The auxiliary coefficients regulate endpoint
alignment, context-field magnitude, context-parameter sparsity, and retention
of the Stage A self-field solution.

\subsection{Numerical trajectory integration}
\label{appb-sec:methods-training-integration}

The endpoint losses and final StageBridge estimand require the solution of the
initial-value problem defined in
Equation~\ref{eq:methods-stagebridge-complete}. Because the neural differential
field has no closed-form integral, its trajectory was approximated
numerically.

The evaluated implementation used the classical fourth-order Runge--Kutta
method, abbreviated RK4. RK4 advances a state across one progression interval
by evaluating the vector field four times: at the beginning of the interval,
at two estimated midpoint states, and at an estimated endpoint state. These
four velocity estimates are combined into a weighted average. This procedure
provides a more accurate approximation of a smooth trajectory than using only
the velocity evaluated at the beginning of each interval.

The interval \([0,1]\) was divided into
\(K_{\mathrm{int}}=4\) equal integration steps. The step width was

\begin{equation}
\delta\tau
=
\frac{1}{K_{\mathrm{int}}}
=
\frac{1}{4}.
\label{appb-eq:methods-rk4-step-width}
\end{equation}

Let \(n\in\{0,\ldots,K_{\mathrm{int}}-1\}\) index an integration step,
let \(\tau_n=n\delta\tau\), and let \(z_n\) denote the numerical approximation
to the regulatory state at \(\tau_n\). For a selected field
\(v(z,w,\tau)\), the four RK4 velocity evaluations were

\begin{equation}
\begin{aligned}
k_{1,n}
&=
v\left(z_n,w,\tau_n\right),
\\
k_{2,n}
&=
v\left(
z_n+\frac{\delta\tau}{2}k_{1,n},
w,
\tau_n+\frac{\delta\tau}{2}
\right),
\\
k_{3,n}
&=
v\left(
z_n+\frac{\delta\tau}{2}k_{2,n},
w,
\tau_n+\frac{\delta\tau}{2}
\right),
\\
k_{4,n}
&=
v\left(
z_n+\delta\tau k_{3,n},
w,
\tau_n+\delta\tau
\right).
\end{aligned}
\label{appb-eq:methods-rk4-slopes}
\end{equation}

Each vector
\(k_{a,n}\in\mathbb{R}^{d_z}\), for \(a\in\{1,2,3,4\}\), is a regulatory
velocity estimate. The first uses the current state. The second and third use
two midpoint state estimates. The fourth uses an estimated state at the end of
the interval.

The state update was

\begin{equation}
z_{n+1}
=
z_n
+
\frac{\delta\tau}{6}
\left(
k_{1,n}
+
2k_{2,n}
+
2k_{3,n}
+
k_{4,n}
\right).
\label{appb-eq:methods-rk4-update}
\end{equation}

The coefficients \(1,2,2,1\) form the standard RK4 weighted average of the
four velocity evaluations. Starting from
\(z_0=z_i^{\mathrm{obs}}\), four applications of
Equation~\ref{appb-eq:methods-rk4-update} produce the numerical endpoint
\(z_4\), which approximates \(z_i(1)\). During full-field integration, the
source context \(w_i\) was supplied unchanged to every field evaluation in
Equation~\ref{appb-eq:methods-rk4-slopes}. The implementation therefore preserved the source-fixed context convention within every intermediate RK4 calculation.


\subsection{Optimization and reproducibility}
\label{appb-sec:methods-training-optimization}

Model parameters were updated with Adam, an adaptive first-order gradient
optimizer. At each step, Adam uses the current gradient together with running estimates of its first and second moments to determine a separate update scale for each parameter.

Stage A used \(1{,}000\) optimization steps with learning rate \(10^{-3}\).
Stage B used \(500\) steps with learning rate \(10^{-3}\). Stage C used
\(500\) steps with learning rate \(10^{-4}\). The learning rate controls the
overall magnitude of each parameter update. Each optimization step used
\(B=64\) transport-sampled pairs and a newly estimated
conditional-stratified transport plan constructed from at most
\(B_{\mathrm{OT}}=512\) source receivers and \(512\) target receivers.

The Euclidean norm of the complete gradient vector was clipped at \(10\).
When the unscaled gradient norm exceeded this threshold, all gradients were
rescaled by a common factor before the optimizer update. This procedure
limited unstable parameter changes arising from an unusually large loss
gradient.

A sampled minibatch occasionally contained no ecological stratum satisfying
the cross-stage common-support requirements. In that case, the minibatch was resampled. A maximum of \(20\) resampling attempts was allowed for each
optimization step. Failure to obtain a valid stratum after these attempts
terminated the fold-specific run and recorded insufficient ecological support for the configured sampling procedure.

Random-number generators used for donor sampling, receiver sampling,
progression-coordinate sampling, and parameter initialization were seeded for each fold. Training and held-out donor sets were checked for disjointness before model fitting. The LUAD analysis used five donor-held-out folds, and the PanIN analysis used three. The fitted parameters, Stage A self-field anchor, training history, maximum coupling dimensions, donor assignments, and frozen coarse-niche transformation were retained for held-out evaluation.


\section{Endpoint metric definitions}
\label{sec:app-b-metrics}

Section~\ref{sec:methods-heldout-endpoints} defines the primary conditional
endpoint metric (Equation~\ref{eq:methods-conditional-endpoint-cost}) and names
three unstratified sensitivity summaries. Those three are defined here.

\subsection*{Global transport cost}

The global metric applies the same entropically regularized squared-Euclidean
transport cost to the unstratified held-out populations,
$D_{\mathrm{OT},f}^{\mathrm{global}} = \operatorname{Sinkhorn}_{\varepsilon}
(\widehat{Z}_{1,f}, Z_{\mathrm{tgt},f})$,
with $\varepsilon = 0.05$ and $80$ Sinkhorn iterations. Because the transport
matrix grows with the product of the two population sizes, each population was
deterministically subsampled to at most $2{,}048$ receivers before this
calculation when required.

\subsection*{Energy distance}

For random vectors $X$ and $Y$ drawn from the predicted-endpoint and
target-stage populations, the population energy distance is

\begin{equation}
D_{\mathrm{E}}(X,Y)
=
\mathbb{E}\!\left[\lVert X-Y \rVert_2\right]
-
\tfrac{1}{2}\,\mathbb{E}\!\left[\lVert X-X' \rVert_2\right]
-
\tfrac{1}{2}\,\mathbb{E}\!\left[\lVert Y-Y' \rVert_2\right],
\label{eq:app-b-energy-distance}
\end{equation}

where $X'$ and $Y'$ are independent draws from the same populations as $X$ and
$Y$ respectively. The empirical estimate therefore compares the mean
between-population distance against the internal dispersion of each population,
and is zero only when the two distributions coincide. Pairwise calculations were
deterministically subsampled to at most $1{,}000$ receivers per population.

\subsection*{Target-stage classifier probability}

Within each held-out fold, a logistic classifier was fitted to the observed
source-stage states as class $0$ and the observed target-stage states as class
$1$, then applied to the predicted endpoints $\widehat{Z}_{1,f}$. The reported
quantity is the mean probability assigned to class $1$, which measures whether
predicted endpoints resemble the observed target stage more closely than the
observed source stage. This score is descriptive: it did not enter model
fitting, model selection, or any falsification-ladder comparison.
